\section{Arithmetic graphs and simultaneous addressing}\label{gr:section}

Fix $b\ge1$.  An \emph{address} is a point $x\in\mathbf F_2^b$;
a vertex with activation $\lambda\in(\mathbf F_2^b)^*$ is present at $x$
when $\lambda(x)=1$.  We will adjoin fresh primes with nonzero activations
to a given collection of primes.  The goal is to make one or two binary
coefficient tests equal to $1$ at every $x\ne0$, while retaining all the
original primes and their residue data.  No new prime is present at $0$.
The hypothesis is that each test, separately at each nonzero address,
is a nonzero polynomial in the available residue bits.

The arithmetic identities are those of
Proposition~\ref{co:symbol-interface}, conditional on its coefficient
lower bounds.  We first describe and realize their graph data, then state
the simultaneous-testing result.  Its proof uses contraction operators
and two finite networks that isolate a prescribed derivative at one address.

\subsection{Labels and arithmetic realization}

Fix a finite set $S$ of rational primes, containing $2$ and all the ramified
primes in the fixed data.  Let $F$ be either $\mathbf Q$ or a quadratic field.
For the rational two-torsion application $F=\mathbf Q(E[2])$; the construction
with full rational two-torsion also uses a quadratic auxiliary field in place
of $F$.  Include in $S$ primes whose prime ideals generate the class group of
$F$, and include the infinite places in $S_F$.  Put
\[
 V_S(F)=\{\alpha\in F^\times/F^{\times2}:
           v_{\mathfrak p}(\alpha)\equiv0\pmod2
           \text{ for }\mathfrak p\notin S_F\}.
\]
This is a finite-dimensional $\mathbf F_2$-space: the valuation sequence
expresses it as an extension of the finite squareclass group of $S_F$-units
by a subgroup of the $2$-torsion of the $S_F$-class group.  Its Kummer extension
is denoted $F_S^{(2)}=F(\sqrt{V_S(F)})$.

A \emph{split vertex} is a rational prime $p\notin S$ split in $F$, together
with a choice $\mathfrak p\mid p$ if $F$ is quadratic.  This choice is its
\emph{orientation}.  Choose a generator $\pi_{\mathfrak p}$ with divisor
$\mathfrak p$ outside $S_F$; class-group generation makes this possible.
Conjugation supplies the generator for the other orientation, and
\begin{equation}\label{gr:norm}
 N_{F/\mathbf Q}(\pi_{\mathfrak p})=p u_{\mathfrak p},
 \qquad u_{\mathfrak p}\text{ a signed rational }S\text{-unit}.
\end{equation}
For $F=\mathbf Q$ use the rational generator $p$.  A \emph{simple vertex}
is an inert rational prime in the quadratic case, with rational generator
$p$.  The fixed label of a split vertex is its Frobenius on $F_S^{(2)}$;
simple vertices have rational fixed-support labels.  A \emph{type} may
additionally record the generators' local squareclasses, valuation parities,
real signs, and prescribed finite ray conditions at $S_F$.  Only compatible,
individually realizable types will be prescribed.  A \emph{filter} fixes the
total rational labels and sign of the signed-prime product
$\prod_p p^*$, where $p^*=(-1/p)p$.

For a pair of split vertices in the quadratic case there are two independent
forward bits: the residues of $\pi_{\mathfrak p}$ and
$\overline{\pi_{\mathfrak p}}$ at the chosen place of the other vertex.
Write quadratic residues additively in $\mathbf F_2$.  Reverse bits are
determined by Hilbert reciprocity and the fixed types.  Rational norm bits
are determined by summing the two oriented bits and the fixed correction
in \eqref{gr:norm}.  When at least one endpoint is simple, or $F=\mathbf Q$,
there is one rational forward bit.  Self-residues are never graph data.
An \emph{arithmetic graph} is this collection of vertices, types, and bits.
A split vertex is \emph{primary} when its generator is a square unit at
every finite place of $S_F$ and positive at all real places.
Different symbols may read compatible projections of the same graph, for
example its rational norm graph.  Here compatibility means that the
projection forgets fixed-label coordinates or takes rational norm data,
carries primary vertices to primary vertices, and commutes with the
omitted-vertex product substitution.  For the norm projection a single
oriented-bit toggle changes the rational bit once.  We use these arithmetic
projections throughout.

\begin{lemma}[Finite realization and product labels]\label{gr:realization}
Fix finitely many old vertices, with all their mutual data, and
individually realizable new types.  Every prescription of the independent
forward bits involving at least one new vertex has a realization by
distinct fresh primes, avoiding an arbitrary finite set.  The old vertices
and their mutual data remain unchanged.  The same holds
with further fixed local tests included in $S$.  A contracted split cluster
may be represented by the product of its oriented generators, with
conjugations prescribed along its contraction tree; its support types and
external residue bits are the corresponding sums.  A cluster containing
one simple vertex has the corresponding rational product data.

For a symbol satisfying the omitted-vertex dependence in
Proposition~\ref{co:symbol-interface}, its value depends only on the graph
data.  Products in the preceding assertion may therefore be replaced by
fresh prime representatives without changing the prescribed contracted
symbol.
\end{lemma}

\begin{proof}
First note that the listed generators account for all the relevant radicals.
If $U$ is a finite set of new primes and a squareclass has even valuations
outside $S_F\cup U_F$, divide by the generators at the primes of $U_F$
where its valuation is odd.  The result belongs to $V_S(F)$.  At an inert
rational prime its rational generator has valuation one at the unique place
of $F$.  Thus fixed-support labels and the displayed residue bits determine
every Frobenius test in an extension with this allowed ramification.

We describe the prime selection carefully.  For a new split vertex take a
template $\pi_0$ of its prescribed type, with prime ideal
$\mathfrak p_0$ away from the old vertices.  Weak approximation gives
$a\in F^\times$ which is sufficiently close to $1$ at $S_F$, positive at
the prescribed real places, and has any prescribed nonzero residues at the
old prime ideals.  Choose those residues so that $a\pi_0$ has the desired
quadratic residues there.  Use the ray modulus consisting of sufficiently
deep powers of the primes of $S_F$, all old prime ideals, and the real
places.  Choose a prime ideal $\mathfrak q$ in the ray class of
$(a)\mathfrak p_0$.  Then
\[
 \mathfrak q=(a\beta)\mathfrak p_0,
 \qquad \beta\equiv1\pmod{\text{ray modulus}},
\]
and $\pi_{\mathfrak q}=a\beta\pi_0$ is the required generator outside $S_F$.
There are infinitely many such degree-one prime ideals.  Indeed, the
prime-ideal theorem in a ray class gives order $X/\log X$ prime ideals of
norm at most $X$, whereas prime ideals of residue degree greater than one
contribute only $O(X^{1/2})$ in a quadratic field.  Degree one gives a split
rational prime, and finite exclusions remove only finitely many choices.

At an old inert prime with residue field $\mathbf F_{p^2}$, prescribe the
norm of the nonzero residue.  The norm onto $\mathbf F_p^\times$ is
surjective, so either rational residue bit is possible.  For a new simple
vertex the requirements are rational congruences, together with the inert
class in $F$.  The Chinese remainder theorem and Dirichlet's theorem supply
fresh primes.  Successively adjoining vertices proves simultaneous finite
realization.  At every step reciprocity determines the reverse bits:
the local Hilbert symbols at the support are fixed, and the only other
contributions are the two endpoint residues.  Taking norms gives exactly
\eqref{gr:norm}, so it imposes no additional independent bit.

The same argument applies to a product template.  Its valuations and local
squareclasses at $S_F$ are those of the product, and its residues at the
remaining vertices multiply.  The ray-class construction replaces that
template by a degree-one prime with those data.  In a simple cluster we
instead use rational congruences for the norm labels.  No residue at a prime
inside the cluster is imposed on its replacement.

In particular this preserves the rational filter, including its norm
correction.  Write $\pi=\prod_i\pi_i$ for the oriented product template
and $\pi'=a\beta\pi$ for its fresh representative; $a\beta$ is a square unit at
$S_F$ and positive at its real places.  The support valuations of $\pi'$
and $\pi$ agree exactly.  Hence
\[
 N\pi=(\prod_i p_i)u,\qquad N\pi'=q u
\]
with the same signed rational $S$-unit $u$: all its valuations and its
sign agree.  The quotient $q/\prod_i p_i=N(a\beta)$ is a local square at
the rational support.  Its squareclass at $2$ also identifies
$q^*$ with $\prod_i p_i^*$ at those fixed tests.  Thus no free choice
of an untracked $S$-unit has entered a product replacement.

Finally, compare two disjoint realizations with identical graph data.
Introduce fresh bridge vertices with matching bits to both realizations,
and replace vertices one at a time.  At each replacement the omitted-vertex
elementary Frobenius tests agree, by the generator description in the first
paragraph.  The symbol interface gives equal values.  A preliminary bridge
removes any overlap between the original realizations.  The same reasoning
applies to the product replacements.
\end{proof}

Primary types are realized by the preceding construction, starting with the principal
ray class.  Reciprocity makes their fixed radical labels trivial.  Their
signed rational contributions are squares at the fixed rational tests.
If stationary vertices are also to have neutral incidences with an
auxiliary, include their primes among these local tests.  In particular,
including the support of a fixed quadratic discriminant $k$ ensures that
every such primary auxiliary splits in $\mathbf Q(\sqrt{k})$.

\subsection{The simultaneous-testing statement}

Separate the given vertices into \emph{terminals} and, if present,
stationary external vertices.  Terminals may be always present or may
have nonzero linear activations; external vertices have fixed types and
incidences throughout.  At each nonzero address the symbol
is a function of the independent bits between the active terminals, with
those types and external data fixed.  We call its unique multilinear
polynomial over $\mathbf F_2$ the \emph{terminal polynomial}.  Its variables
describe formal changes of the terminal bits, while the prescribed prime
realization supplies a fixed evaluation point.

For a bit $z$ define the finite difference
$\Delta_z f=f|_{z=0}+f|_{z=1}$.  On a multilinear polynomial it is the
formal partial derivative.  In particular, differences commute and
$\Delta_z^2=0$.

\begin{definition}[Witnesses and permitted partitions]\label{gr:partitionable}
A \emph{witness} is a finite graph on which the
symbol is $1$.  A graph is \emph{partitionable} for the symbol if its vertices
can be partitioned into nonempty clusters, each with at most one simple
vertex, with a choice of oriented contraction trees whose contracted graph
is a witness.
This definition includes a specified contraction certificate; it makes no
assertion about a Selmer group.
\end{definition}

\begin{theorem}[Address lemma]\label{gr:address}
Suppose the binary symbols under consideration depend only on the
arithmetic graph data and satisfy the contraction, whole-rational-prime
omission, and primary-isolation rules of
Proposition~\ref{co:symbol-interface}, with the realization freedoms
of Lemma~\ref{gr:realization}.  Contraction includes the product labels
and external incidences, and the zero rule for two simple vertices.
These identities are required only on nonempty configurations and for
contractions with nonempty output.
Fix $b\ge1$ and either of the following terminal arrangements.
\begin{enumerate}
\item One symbol, a nonempty set of always active terminals, and a
nonzero polynomial in their independent pair bits.
\item One or two symbols on terminals with nonzero linear activations,
including at least one terminal for every nonzero linear form.  At each
$x\ne0$ each symbol separately is a nonzero polynomial in the active
terminal bits.  The symbols may read different compatible projections
of a common graph alphabet and may have fixed stationary external data.
\end{enumerate}
Suppose also that at each nonzero address the total active terminal data,
together with any stationary external data, obey the required filter.
Starting from any fixed realization of the terminal and stationary
vertices, one can adjoin finitely many fresh primary auxiliary vertices,
activated by nonzero linear forms, and choose only incidences involving
these new vertices so that every symbol equals $1$ at every $x\ne0$.
All data on the original vertices remain unchanged.  No auxiliary is
present at $0$.  All auxiliaries may satisfy the fixed splitting and
neutral external-incidence conditions incorporated into the primary type.

It is enough that a terminal polynomial have a nonzero evaluation or
a nonzero iterated edge difference.  In particular a partitionability
certificate suffices for the stationary hypothesis.
\end{theorem}

The proof has two algebraic steps.  Forests of auxiliary vertices act as
finite-difference operators on the original terminal polynomial.  We then
build networks whose degree-one terms select a chosen derivative at a
single address.  Applying them to top-degree coefficients will give the
simultaneous value $1$, independently of the fixed terminal bits.

\subsection{From arithmetic contraction to forest operators}

\begin{lemma}[Oriented forests]\label{gr:forest}
Assume the graph-symbol rules stated in Theorem~\ref{gr:address}.
Differentiation in a tree of graph
bits contracts the tree, adding product labels and external incidences;
its value is zero if the tree contains two simple vertices.  Differentiation
in any collection of edge variables containing a cycle is zero, even when
different edges use incompatible relative orientations.  Two different
oriented variables on the same pair count as a cycle of length two.
\end{lemma}

\begin{proof}
The one-edge assertion is precisely the arithmetic contraction rule,
with product replacements justified by Lemma~\ref{gr:realization}.
Induction proves the tree assertion: a split--split merger is split, a
split--simple merger is simple, and the first merger of two components
containing simple vertices has zero derivative.  Along a tree there is
a unique choice of relative orientations once one root orientation is
fixed.  Its product template depends on support types and external bits,
not on the internal bits being differentiated.

For clarity, consider first the two variables on a split pair $p,q$.
Differentiation in the inertia bit at one place over $p$ gives the symbol
formed with the \emph{entire rational prime $p$ omitted}.  Its elementary
extension is unramified at both places over $p$.  Differentiating in the
other oriented bit is therefore zero.  Repetition of the same variable
is zero by $\Delta_z^2=0$.  This argument would not follow from omission
of only one oriented prime ideal.

For a longer cycle, first contract a spanning path.  By the product-label
rule the closing variable has become internal to the merged cluster.
Its change affects neither the product's support data nor any incidence
to a remaining vertex.  Equivalently it is an inertia change at an omitted
endpoint, covered by the preceding unramifiedness argument.  Thus its
derivative vanishes.  Possible conjugation of the product generator only
changes the induced orientation and does not reintroduce that endpoint.
Since there is no self-residue among the inputs to the replacement symbol,
there is no extra closing-edge datum to retain.  Additional derivatives
commute with this zero derivative.

For example, contract the aligned edge $12$ in a triangle.  With
$\pi_{12}=\pi_1\pi_2$, the two bits toward vertex $3$ are
$t_0=z_{13}^0+z_{23}^0$ and $t_1=z_{13}^1+z_{23}^1$, up to fixed
type corrections.  An aligned edge $23$ then differentiates $t_0$;
a conjugately aligned closing edge $31$ differentiates $t_1$.
Their mixed derivative is the omitted-endpoint zero just proved.
If the closing edge is aligned instead, both become the same
difference and its square is zero.  Contracting a longer spanning
path reduces to exactly these two alternatives.
\end{proof}

Fix a realizable baseline terminal graph and let $P$ be its terminal
polynomial.  A \emph{network} consists of additional primary
vertices and programmed edges, with no direct terminal--terminal edges.
All unprogrammed incidences of its primary vertices are neutral.  Within
one network choose one orientation at each terminal and consistently use
the corresponding edge alignment.  Networks may choose different
alignments.  The \emph{degree} of a contraction is the number of terminal
vertices it removes by merging them.

\begin{lemma}[Network expansion and its first two degrees]\label{gr:network}
A network acts on $P$ by a sum of finite-difference operators.  A surviving
term is a forest in which every component containing a primary vertex
meets a terminal.  A component meeting $t$ terminals acts by a contraction
of degree $t-1$; a component containing two simple terminals is zero.
Networks on disjoint sets of internal vertices, with no edges between
those sets, multiply as operators on the same terminal polynomial.
Consequently an operator of degree greater than $\deg P$ annihilates $P$.

Let $Q$ be the grounded Laplacian on the internal vertices: $Q_{ij}$ for
$i\ne j$ records the internal edge $ij$, while $Q_{ii}$ is the total
incidence at $i$, including its terminal edges.  Let $(c_a)_i$ record the
incidence from internal vertex $i$ to terminal $a$.  All entries are in
$\mathbf F_2$.  The constant and pair-contraction coefficients of the
operator are respectively
\begin{equation}\label{gr:forest-coefficients}
 \det Q,\qquad c_a^{\mathsf t}\operatorname{adj}(Q)c_b\quad(a\ne b).
\end{equation}
The second coefficient multiplies the terminal difference in the alignment
selected by the network.  A prohibited pair multiplies the zero operator.
\end{lemma}

\begin{proof}
Expand the evaluation at the programmed incidences from the graph with
all primary incidences zero.  For a set $E_0$ of programmed bits this is
the elementary identity
\[
 f(\mathbf1)=\sum_{T\subseteq E_0}(\Delta_T f)(\mathbf0).
\]
Lemma~\ref{gr:forest} kills cyclic $T$.  A tree component meeting no
terminal contracts to an isolated primary vertex, and the isolation rule
kills it.  This includes an untouched primary vertex.  In any remaining
component all primary labels are trivial and all external primary bits
at the expansion baseline are zero.  Its contracted label and its
incidences are exactly those obtained by contracting its terminal
vertices in the original graph.  Choose any terminal tree inducing its
relative orientations; the corresponding terminal differences therefore
give the same value.  Its degree is the number of lost terminal vertices,
namely $t-1$.  This establishes the operator interpretation, rather than
merely a count of possible contraction degrees.

Expanding two disjoint networks at once gives the composition of these
operators: the product-label rule gives the same final labels in either
order, and any cycles created between networks vanish by
Lemma~\ref{gr:forest}.  Ordinary polynomial differentiation now proves
the asserted degree bound.

This argument also explains the compatible projections.  Taking norms
commutes with every oriented product in a forest; primary factors have
neutral norm corrections, and each programmed oriented toggle projects
to its one rational toggle.  Thus the projected forest acts by the
corresponding differences on the projected \emph{original} terminal
polynomial at every degree.  If projected differences repeat an edge
or merge two simple clusters, that term is zero.  Forgetting fixed-label
coordinates plainly commutes with the same products.  Each projected
network therefore acts on the original projected terminal polynomial.

The determinant counts are instances of the all-minors matrix-tree
theorem \cite[\S1]{Chaiken1982}.  We give the Cauchy--Binet proof
in the form needed here, including singular $Q$.  Let $I$
be the internal vertex set and $A$ the terminal set.  Write the full
Laplacian as $L=DD^{\mathsf t}$, where a column of the incidence matrix
$D$ has ones at its edge's endpoints.  Cauchy--Binet applied to
$L_{I,I}=Q$ counts edge sets whose incidence matrix with the terminal
rows deleted is invertible.  Such an edge set is a forest with exactly
one terminal in each component: a cycle gives dependent columns; a
component with no deleted row gives dependent rows; and for a tree with
one row deleted, leaf removal proves invertibility.  These are exactly
the degree-zero terms.

For the pair coefficient apply Cauchy--Binet with row set $I\cup\{a\}$
and column set $I\cup\{b\}$.  Each counted forest has one root from
$A\setminus\{a\}$ and one from $A\setminus\{b\}$ in each component.
This requires that $a,b$ lie together and that every other terminal lie
in its own component.  The corresponding Laplacian minor is
\[
 \det\begin{pmatrix}Q&c_b\\c_a^{\mathsf t}&0\end{pmatrix}
       =c_a^{\mathsf t}\operatorname{adj}(Q)c_b.
\]
The corner is zero because the network has no direct terminal edge;
signs disappear in characteristic two.  The determinant identity uses
the adjugate and so remains valid when $Q$ is singular.
\end{proof}

\subsection{Two networks which select one address}

Fix a nonzero target $x_*$ and a terminal bit $z$ joining two terminals
$a_0,a_1$ active at $x_*$.  In the
moving case assume that some terminal has each nonzero activation form;
in the stationary case all terminals are always present.

Choose a basis $\lambda_1,\ldots,\lambda_b$ of the dual space with
$\lambda_i(x_*)=1$ for every $i$.  To obtain it, choose $\lambda_1(x_*)=1$
and replace a basis $\mu_2,\ldots,\mu_b$ of the annihilator of $x_*$ by
$\lambda_1+\mu_2,\ldots,\lambda_1+\mu_b$.
Make $b$ groups of three primary vertices, group $i$ activated by
$\lambda_i$.  Since these forms are a basis, all groups are present
exactly at $x_*$.  At every other address at least one group is absent.
In each physical three-dimensional vertex space choose
a basis $(e_i,u_i,v_i)$ with $e_i=(1,1,1)$.  Thus the physical coordinate
change is block diagonal by groups.

Prescribe a symmetric bilinear form $Q$ by making the $e_i$ orthonormal
and orthogonal to all $u_j,v_j$, and putting on the remaining coordinates
\begin{equation}\label{gr:H}
 H(T,B)=\begin{pmatrix}0&T\\T^{\mathsf t}&B\end{pmatrix}.
\end{equation}
Matrix entries in the physical coordinates are denoted by the same
letter $Q$.  For a subset $J$ of active groups, $Q_J$ means its physical
principal restriction.  Define the column $h_i=Qe_i$.
It is supported on group $i$: in adapted coordinates it is the covector
dual to $e_i$, and the coordinate change does not mix groups.  Its value
is independent of $T$ and $B$.

Attach $h_i$ to an anchor terminal with activation $\lambda_i$.  In the
stationary case attach it to any fixed terminal instead.  The support
property gives, on every address,
\[
 \sum_{i\in J}h_i|_J=Q_J\mathbf1_J.
\]
Program internal edges using the off-diagonal
physical entries of $Q$.  A matrix is their grounded Laplacian precisely
when the sum of the terminal incidence columns is its row-sum vector,
so these anchor columns give exactly $Q_J$, including its diagonal.

We shall also attach the same additional column $c$ to $a_0$ and $a_1$.
In the moving case with distinct activation forms $\alpha,\beta$, attach
$c$ also to a terminal of activation $\alpha+\beta$; this terminal is
inactive at $x_*$.  If $\alpha=\beta$, no third column is needed.  The
sum of the additional active columns is zero at every address, since
$\alpha+\beta+(\alpha+\beta)=0$.  In the stationary case the two columns
cancel directly.  Coincident terminal roles mean addition of their
columns.  This preserves the prescribed grounded matrix after every
restriction.

For a concrete instance, take $b=1$, target $x_*=1$, and two always-active
split terminals $a_0,a_1$ with selected bit $z$.  Label the three primary
vertices $1,2,3$, all activated by $x$.  Choose
$e=(1,1,1)^{\mathsf t}$, $u=(1,0,0)^{\mathsf t}$,
$v=(0,1,0)^{\mathsf t}$, and $T=(0)$, $B=(1)$.
The form in this adapted basis is $\operatorname{diag}(1,0,1)$.
In physical coordinates $h=Qe=(0,0,1)^{\mathsf t}$ and the covector
dual to $u$ is $c=(1,0,1)^{\mathsf t}$.  Using $a_0$ as the anchor gives
\[
 Q=\begin{pmatrix}0&0&0\\0&1&1\\0&1&0\end{pmatrix},\qquad
 c_{a_0}=h+c=\begin{pmatrix}1\\0\\0\end{pmatrix},\qquad
 c_{a_1}=c=\begin{pmatrix}1\\0\\1\end{pmatrix}.
\]
Thus the programmed graph is the path
$a_0\! -\! 1\! -\! a_1\! -\! 3\! -\! 2$.
Here $\det Q=0$ and $\operatorname{adj}(Q)=uu^{\mathsf t}$, so
$c_{a_0}^{\mathsf t}\operatorname{adj}(Q)c_{a_1}=1$.
There are only two terminals, so no higher contraction degree occurs:
the network acts exactly by $\Delta_z$ at $1$.  At $0$ all three primary
vertices are absent, and the network acts by the identity.

Write $O_{\ge d}$ for a sum of contraction operators of degrees at least
$d$; it is a filtration notation, not an estimate.

\begin{lemma}[The two address programs]\label{gr:programs}
With the preceding grounding convention one has the following programs
for every $b\ge1$.
\begin{enumerate}
\item A network $\mathcal S$ with
\[
 \mathcal S_{x_*}=\Delta_z+O_{\ge2},\qquad
 \mathcal S_x=1+O_{\ge1}\quad(x\ne x_*).
\]
\item Two networks on the same activated sites, indexed by
$\varepsilon\in\mathbf F_2$, with
\[
 \mathcal N_x^{(\varepsilon)}=1+O_{\ge1},\qquad
 \mathcal N_x^{(1)}+\mathcal N_x^{(0)}
 =\begin{cases}\Delta_z+O_{\ge2},&x=x_*,\\
 O_{\ge2},&x\ne x_*.
 \end{cases}
\]
\end{enumerate}
All nonzero-activation sites are absent at $0$.  Unspecified incidences,
including those to stationary external vertices, may remain neutral.
\end{lemma}

\begin{proof}
For the first program take $T=I+P$, where $P$ is the cyclic permutation
matrix on $b$ coordinates, and $B$ diagonal with exactly one nonzero
entry.  A proper principal group restriction of $P$ is a union of
directed paths, hence nilpotent.  Thus $T_J$ is invertible for every
proper $J$, and so is $H(T_J,B_J)$: its kernel equations first give
$v=0$ and then $u=0$.  This includes the empty restriction, whose
determinant is $1$.

On the full group set, both $\ker T$ and $\ker T^{\mathsf t}$ are
spanned by the all-one vector.  A radical vector of $H$ has
$v=t\mathbf1$; summing the equation $T^{\mathsf t}u+Bv=0$ gives $t=0$
because $\mathbf1^{\mathsf t}B\mathbf1=1$.  Thus the full $Q$ has
one-dimensional radical spanned by $r=\sum_i u_i$.  In characteristic
two its adjugate is $rr^{\mathsf t}$.  Indeed a rank-$n-1$ matrix has a
nonzero rank-one adjugate with image its kernel; symmetry identifies
the two factors, and the only nonzero scalar is $1$.
Choose $c$ to be the covector dual to $u_1$, so $c^{\mathsf t}r=1$.
All anchors pair to zero with $r$.  At $x_*$ the only columns pairing
nontrivially with $r$ are consequently those of $a_0,a_1$.  Formula
\eqref{gr:forest-coefficients} gives zero constant and exactly
$\Delta_z$ in degree one.  At every other address the group subset is
proper and the constant is $1$.  When $b=1$, $I+P=0$ and $B=(1)$,
so the same radical calculation applies without an exception.

For the second program let $T=I+L$, where $L$ is the consecutive lower
shift, and take $B_\varepsilon=\varepsilon E_{bb}$.  Every group
restriction of $T$ is invertible, and block multiplication gives
\begin{equation}\label{gr:inverse}
 H(T,B)^{-1}=
 \begin{pmatrix}T^{-\mathsf t}BT^{-1}&T^{-\mathsf t}\\
 T^{-1}&0\end{pmatrix}.
\end{equation}
Again take $c$ dual to $u_1$.  On a group subset $J$, the change of
$c^{\mathsf t}Q_J^{-1}c$ is
\[
 \bigl((T_J^{-1})_{b1}\bigr)^2,
\]
interpreted as zero if $1$ or $b$ is absent.  Since
$(I+L_J)^{-1}=I+L_J+L_J^2+\cdots$, that entry is $1$ precisely when
the entire path $1,2,\ldots,b$ is present.  This is exactly the target
address.  Pairings involving an anchor do not change:
$Q_J^{-1}h_i=e_i$, and the vectors $e_i,h_i,c$ do not depend on
$\varepsilon$.  At the target the balancing terminal is absent, so
the only changed pair coefficient is the one for $a_0,a_1$.
Every constant is $1$.  Lemma~\ref{gr:network} proves the claims,
including $b=1$, when the path consists of its single vertex.
\end{proof}

\subsection{Proof of the address lemma}

\begin{proof}[Proof of Theorem~\ref{gr:address}]
Use the prescribed realization of the terminal and stationary vertices
as a common baseline for all terminal bits.  Regard
their toggles as common variables, omitting variables invisible to a
given symbol.  A nonzero iterated difference implies that the polynomial
is nonzero.  Conversely, the unique multilinear polynomial of a
nonzero function has a top monomial with nonzero coefficient.  These
observations also explain the last assertion using
Lemma~\ref{gr:forest}.
Every nonzero address has an active terminal, and contractions never
remove the last vertex.  Thus only nonempty configurations are used.

First consider one symbol.  At each $x\ne0$ let $P_x$ have degree $d_x$
and choose a monomial with coefficient $1$ in its top homogeneous part.
For each variable of this monomial add a separate singular program
$\mathcal S$ targeted at $x$.  Give all these networks disjoint internal
vertices.  At address $x$, its own $d_x$ networks each have zero constant,
so every surviving term spends at least $d_x$ contraction degrees.
An extra degree from any network kills $P_x$.  Thus only the designated
degree-one contractions of its own networks and the constant terms
of all other networks survive.  Their value is the chosen top
coefficient, namely $1$.  A nonzero constant polynomial needs no
network and already has value $1$.  This argument applies equally to
always active terminals, since Lemma~\ref{gr:programs} supplies their
stationary anchors.

For two symbols write their polynomials at address $x\ne0$ as
$P_x,R_x$, of degrees $d_x,e_x$, and let $P_x^{\rm top},R_x^{\rm top}$
be their top homogeneous parts.  Multiply these parts in the ordinary
polynomial ring
$\mathbf F_2[z_1,\ldots,z_m]$, \emph{before} imposing any Boolean
relations.  That ring is an integral domain, so choose a coefficient
equal to $1$ at a monomial
\begin{equation}\label{gr:chosen-monomial}
 M_x=\prod_z z^{\nu_{xz}},\qquad
 \nu_{xz}\in\{0,1,2\},\qquad
 \sum_z\nu_{xz}=d_x+e_x.
\end{equation}
For $\nu_{xz}=2$ add one singular program for $z$ targeted at $x$.
For $\nu_{xz}=1$ add one nonsingular program for $z$ targeted at $x$,
with its own choice $\varepsilon_{xz}\in\mathbf F_2$.  All internal
vertices of different programs are disjoint.  Let
\[
 \Phi(\varepsilon)=\prod_{x\ne0}
       \bigl(\text{final evaluation of }P_x\bigr)
       \bigl(\text{final evaluation of }R_x\bigr).
\]
We prove that the sum of $\Phi$ over all choices of $\varepsilon$ is $1$.
This proves that some single choice has all the required evaluations
equal to $1$.

Expand every program in contraction degree using
Lemma~\ref{gr:network}.  A singular program targeted at $x$ has zero
constant in both factors there, so it costs at least one degree in
$P_x$ and one in $R_x$.  A nonsingular program has constant $1$
independent of its choice.  In the sum over that choice, a term using
only its constants in all factors cancels.  Every remaining term must
therefore spend at least one degree somewhere for this program.  The
sum of these minimum costs is exactly
\begin{equation}\label{gr:budget}
 D=\sum_{x\ne0}(d_x+e_x),
\end{equation}
by \eqref{gr:chosen-monomial}.  A nonzero term cannot exceed this total:
if it did, some factor would receive more differences than its degree.
Hence every surviving term uses every minimum cost exactly.

It follows that a singular program spends exactly one degree in each
factor at its target and zero elsewhere.  Its degree-one term there
is precisely $\Delta_z$.  Each nonsingular program spends exactly one
degree in exactly one factor.  After summing its choice, that
degree-one contribution is zero away from its target, and is
$\Delta_z$ at its target.  All other occurrences of these programs
contribute their common constant $1$.  Thus no unknown higher forest
coefficient enters the surviving calculation.

At a fixed address, let $Z_2$ be the variables with exponent two and
$Z_1$ those with exponent one.  The local allocation sum is
\[
 \sum_{J\subseteq Z_1}
 [z_{Z_2}z_J]P_x^{\rm top}\,
 [z_{Z_2}z_{Z_1\setminus J}]R_x^{\rm top}
   =[M_x](P_x^{\rm top}R_x^{\rm top})=1.
\]
Coefficients of the wrong degree in this display are zero.  Every
factor must in fact receive its full degree, since the total is $D$
and none can exceed its own degree.  Therefore its final derivative
is exactly the displayed top coefficient, independent of the chosen
baseline values.  Allocations at different addresses are independent,
so their product is $1$, as asserted.  This also covers a constant
factor, whose degree is zero.  For example, when both top forms are
$z$, their ordinary product is $z^2$ and the singular program
differentiates each factor once; Boolean reduction would lose this
allocation.

In the two-symbol case, choose the program parameters for which
$\Phi=1$.  In either case, apply Lemma~\ref{gr:realization} to adjoin
the resulting finite networks to
the fixed original vertices.  Only new vertices and incidences involving
them have been prescribed.  The primary types preserve the filters.
All activations are nonzero linear forms,
so no new vertex appears at $0$, and the stipulated neutral and
splitting conditions persist.
\end{proof}

The theorem is finite for each $b$; it does not assert a single infinite
array of primes.  For a stabilized symbol, first choose $b$, its finite
terminal data, and all the networks above.  Their total weight and the
list of coefficient tests are then finite.  A sufficiently advanced
stage at which the symbol interface holds for that weight range
realizes the construction.  Required precision and the number of
auxiliaries may grow with $b$.  The arithmetic valuation constant
attached to a unit symbol must still be supplied by the coefficient
construction independently of $b$; no determinant or forest count in
this section proves that uniformity.
